\chapter{From binomial to Black--Scholes}
\label{ch:black-scholes}

\begin{question}
What happens to the binomial price of a call when each period becomes tiny and the number of periods huge?
\end{question}

\begin{whatfor}
\begin{itemize}
  \item Writes the binomial call price as a difference of two binomial probabilities.
  \item Shrinks the steps ($u_n=e^{\sigma\sqrt{T/n}}$, $d_n=e^{-\sigma\sqrt{T/n}}$, $1+r_n=e^{rT/n}$) and uses the central limit theorem.
  \item Arrives at the Black--Scholes formula $C=S_0\Phi(h)-Ke^{-rT}\Phi(h-\sigma\sqrt T)$.
\end{itemize}
\par\smallskip\textbf{Skipping this chapter?} Its formulas are rows of the continuous-time pricing card, in the formula sheet on page~\pageref{fs:mg-pricing}.
\end{whatfor}

\begin{prereq}
From Chapter~21: the $T$-period binomial price and the split of the call into ``in the money'' terms. From Chapters~8 and~11: the central limit theorem and geometric Brownian motion. From Part~0: $\Phi$.
\end{prereq}

%=====================================================================
\section{The call as two binomial probabilities}

In a $T$-period tree the call pays $(u^jd^{T-j}S_0-K)^+$ after $j$ up-moves. It pays only when $j\ge m$, where $m$ is the least $j$ with $u^jd^{T-j}S_0>K$; taking logarithms,
\[
  m=\Big\lfloor\frac{\ln K-\ln S_0-T\ln d}{\ln u-\ln d}\Big\rfloor+1 .
\]
Splitting the two terms of $(u^jd^{T-j}S_0-K)$ in Chapter~21's formula,
\[
  C=S_0\,P(\hat U\ge m)-\frac{K}{(1+r)^T}\,P(U^*\ge m),\qquad U^*\sim B(T,p^*),\quad \hat U\sim B(T,\hat p),\quad \hat p=\frac{up^*}{1+r}. \tag{$\star$}
\]
($\hat p\in(0,1)$ because $up^*+d(1-p^*)=1+r$.)

%=====================================================================
\section{Many small steps}

Divide $[0,T]$ into $n$ steps of length $\delta_n=T/n$, and let
\[
  u_n=e^{\sigma\sqrt{\delta_n}},\qquad d_n=e^{-\sigma\sqrt{\delta_n}},\qquad 1+r_n=e^{r\delta_n}.
\]
After $n$ steps, $S_T=S_0e^{Z_n}$ with $Z_n=Y_1+\dots+Y_n$, $Y_k=\pm\sigma\sqrt{\delta_n}$. With $e^x\approx1+x$,
\[
  p_n^*=\frac{e^{r\delta_n}-e^{-\sigma\sqrt{\delta_n}}}{e^{\sigma\sqrt{\delta_n}}-e^{-\sigma\sqrt{\delta_n}}}\approx\frac12+\frac{r}{2\sigma}\sqrt{\delta_n}\quad\Big(\text{more precisely }\tfrac12+\tfrac{r-\sigma^2/2}{2\sigma}\sqrt{\delta_n}\Big).
\]
This is the walk of Chapter~8 with drift: under $P^*$, $Z_n$ has mean $n\sigma\sqrt{\delta_n}(2p^*_n-1)\to(r-\frac12\sigma^2)T$ and variance $\to\sigma^2T$, so by the central limit theorem
\[
  Z_n\ \longrightarrow\ N\big((r-\tfrac12\sigma^2)T,\ \sigma^2T\big)\quad\text{under }P^*.
\]
The same argument under $\hat p_n$ gives mean $(r+\frac12\sigma^2)T$ and the same variance.

\begin{theorem}[Black--Scholes formula]
\label{thm:bs}
As $n\to\infty$, the binomial price of a European call converges to
\[
  \boxed{\;C=S_0\,\Phi(h)-Ke^{-rT}\,\Phi(h-\sigma\sqrt T),\qquad h=\frac{\ln(S_0/K)+(r+\frac12\sigma^2)T}{\sigma\sqrt T}.\;}
\]
\end{theorem}
\begin{proof}
Take the limit of each probability in $(\star)$.
\begin{steps}
  \item $U^*\ge m$ means the price ends above $K$: $Z_n>\ln(K/S_0)$. With the normal limit,
  \[
    \begin{aligned}
    P(U^*\ge m)&\to P\Big(N\big((r-\tfrac12\sigma^2)T,\sigma^2T\big)>\ln\tfrac{K}{S_0}\Big)\\
    &=\Phi\Big(\frac{\ln(S_0/K)+(r-\frac12\sigma^2)T}{\sigma\sqrt T}\Big)=\Phi(h-\sigma\sqrt T).
    \end{aligned}
  \]
  \item Likewise $P(\hat U\ge m)\to\Phi(h)$, using the mean $(r+\frac12\sigma^2)T$.
  \item $(1+r_n)^n=e^{rT}$. Substitute into $(\star)$. \qedhere
\end{steps}
\end{proof}

\begin{pic}
\begin{tikzpicture}
\begin{axis}[width=10cm,height=4.8cm,xmin=0,xmax=100,ymin=9.4,ymax=12.3,xlabel={number of steps $n$},ylabel={call price},legend style={font=\scriptsize}]
  \addplot[defc,thick,mark=*,mark size=0.6pt] table[x=n,y=c] {figs/bs_conv.dat}; \addlegendentry{binomial price}
  \addplot[trapc,very thick,dashed,domain=0:100] {10.4506}; \addlegendentry{Black--Scholes $10.45$}
\end{axis}
\end{tikzpicture}
\picnote{$S_0=K=100$, $\sigma=0.2$, $r=0.05$, $T=1$. The binomial price zig-zags (odd and even $n$ alternate) and settles on the Black--Scholes value.}
\end{pic}

\paragraph{The limiting model.} In the limit the stock is $S_t=S_0e^{\mu t+\sigma W_t}$, a geometric Brownian motion, and the bond is $B_t=e^{rt}$. Under the risk-neutral probability the exponent's drift is $r-\frac12\sigma^2$, whatever the real-world drift $\mu$. Chapter~23 derives the same price directly in continuous time.

\begin{example}[A Black--Scholes price]
\label{ex:bs}
$S_0=80$, $K=85$, $T=1$, $r=0.05$, $\sigma=0.15$. Then $h=\frac{\ln(80/85)+0.06125}{0.15}\approx0.0042$, and
\[
  C=80\,\Phi(0.0042)-85e^{-0.05}\Phi(-0.1458)\approx80\times0.5017-80.85\times0.4420\approx4.39 .
\]
\end{example}
